\clearpage
\begin{examplebox}
\textbf{\textcolor{CExample}{例1：已知顶角，先认 $b^2$ 再取半角。}}
双曲线 $x^2/4-y^2/5=1$，焦点顶角为 $60^\circ$。
负号项分母给出 $b^2=5$，故
\[
 S=5\cot30^\circ=\boxed{5\sqrt3}.
\]
回验存在性：$c=3$，取 $P=(4\sqrt6/3,5\sqrt3/3)$，
代回方程得 $8/3-5/3=1$；底高给 $S=3\cdot5\sqrt3/3=5\sqrt3$。

\medskip
\textbf{\textcolor{CExample}{例2：面积反求顶角，最后要把半角乘 $2$。}}
双曲线 $x^2/9-y^2/16=1$，焦点三角形面积为 $16$。
$b^2=16$，所以 $\cot(\theta/2)=16/16=1$。
由 $0<\theta/2<90^\circ$ 得 $\theta/2=45^\circ$，故 $\boxed{\theta=90^\circ}$。
回验：$c=5$，可取 $P=(3\sqrt{41}/5,16/5)$，代回方程得 $41/25-16/25=1$，
底高面积为 $5\cdot16/5=16$。

\medskip
\textbf{\textcolor{CExample}{例3：由焦距、面积和顶角反求方程。}}
横向双曲线的焦距为 $10$，$S=8\sqrt3$、$\theta=60^\circ$。
先求虚半轴平方，再用双曲线参数关系：
\[
 b^2=S\tan30^\circ=8,\quad c=5,\quad a^2=c^2-b^2=17>0.
\]
因此 $\boxed{x^2/17-y^2/8=1}$。例6继续用具体点回验全部条件。

\medskip
\textbf{\textcolor{CExample}{例4（最低门槛：勾股与底高双回验）}}
\textbf{题目：} 在双曲线 $x^2-\frac{y^2}{3}=1$ 上取点 $P(2,3)$，证明该点的焦点三角形面积与半角公式一致。
\textbf{解：} $a^2=1,b^2=3,c^2=4$，两焦点为 $F_1(-2,0),F_2(2,0)$。
\[
PF_1=\sqrt{4^2+3^2}=5,\quad PF_2=3,\quad F_1F_2=4.
\]
\textbf{方法A（底高直接求）：} $S=\frac12\cdot4\cdot3=6$。
\textbf{方法B（夹角公式回验）：}
\[
\cos\theta=\frac{5^2+3^2-4^2}{2\cdot5\cdot3}=\frac35,
\quad \sin\theta=\frac45,
\quad\cot\frac\theta2=\frac{1+3/5}{4/5}=2.
\]
所以 $b^2\cot\frac\theta2=3\times2=6$。
\textbf{思考：} 这说明\emph{焦点连线作底求面积}与\emph{两腰夹角求面积}只是同一个三角形的两种算法。

\medskip
\end{examplebox}
\clearpage
\begin{examplebox}
\textbf{\textcolor{CExample}{例5（直角焦点三角形的快速设点）}}
\textbf{题目：} 在双曲线 $x^2-\frac{y^2}{3}=1$ 上，找一个满足 $\angle F_1PF_2=90^\circ$ 且 $y>0$ 的点，并回验面积。
\textbf{解：} 因为 $\cot45^\circ=1$，面积公式立即给出 $S=b^2=3$。
以焦点连线为底，有 $S=c|y|=2y$，故 $y=\frac32$。代回双曲线方程：
\[
x^2=1+\frac{(3/2)^2}{3}=\frac74,
\quad P=\left(\frac{\sqrt7}{2},\frac32\right)
\quad\text{或}\quad
P=\left(-\frac{\sqrt7}{2},\frac32\right).
\]
取右支上的第一个点。它到两焦点的距离分别为 $\sqrt7+1$ 和 $\sqrt7-1$，故
\[
(\sqrt7+1)^2+(\sqrt7-1)^2=16=(2c)^2.
\]
由勾股定理逆定理，$\angle F_1PF_2=90^\circ$，回验成功。
\begin{center}\begin{tikzpicture}[font=\small,>=stealth]
\draw[black!30,rounded corners=2pt] (0,0) rectangle (14.4,9.2);
\node[anchor=west,font=\small\bfseries] at (.5,8.7) {同一双曲线上：越远离顶点，高越大、面积越大，顶角反而越小};
\foreach \cx/\hh in {2.4/0,7.2/1.5,12/6}{
\begin{scope}[shift={(\cx,3.7)},scale=.4]
\draw[black!35] (-2.5,0)--(4.5,0);
\pgfmathsetmacro{\px}{sqrt(1+\hh*\hh/3)}
\fill[orange!12] (-2,0)--(\px,\hh)--(2,0)--cycle;
\draw[blue,very thick,domain=-1:6.6,samples=101,variable=\t] plot ({sqrt(1+\t*\t/3)},\t);
\draw[blue,very thick,domain=-1:6.6,samples=101,variable=\t] plot ({-sqrt(1+\t*\t/3)},\t);
\draw[orange!85!black,very thick] (-2,0)--(\px,\hh)--(2,0)--cycle;
\draw[violet,dashed,thick] (\px,0)--(\px,\hh);
\fill (-2,0) circle (3pt);\fill (2,0) circle (3pt);
\fill (\px,\hh) circle (3pt) node[above left=6pt] {$P$};
\end{scope}}
\node at (2.4,7.5) {顶点 $P=(1,0)$};
\node at (7.2,7.5) {例5：$P=(\sqrt7/2,3/2)$};
\node at (12,7.5) {较远点 $P=(\sqrt{13},6)$};
\node at (2.4,2.45) {$h=0,\quad S=0$};
\node at (7.2,2.45) {$h=3/2,\quad S=3$};
\node at (12,2.45) {$h=6,\quad S=12$};
\node at (2.4,1.65) {$\theta=180^\circ$（退化）};
\node at (7.2,1.65) {$\theta=90^\circ$};
\node at (12,1.65) {$\theta\approx28.07^\circ$};
\node at (7.2,.6) {同尺度、同曲线 $x^2-y^2/3=1$；$c=2$，$S=2h=3\cot(\theta/2)$。};
\end{tikzpicture}\end{center}

中图正是本例的点；左右两图分别补上顶点与较远位置。

\medskip
\end{examplebox}
\clearpage
\begin{examplebox}
\textbf{\textcolor{CExample}{例6（反推参数时先检查可行性）}}
\textbf{题目：} 已知焦点在 $x$ 轴的双曲线焦距为 $10$，其某焦点三角形满足 $S=8\sqrt3$、$\theta=60^\circ$。求方程并说明所给角度确能实现。
\textbf{解：} $b^2=S\tan30^\circ=8$，$c=5$，所以 $a^2=c^2-b^2=17>0$，方程为
\[
\boxed{\frac{x^2}{17}-\frac{y^2}{8}=1}.
\]
再用 $S=c|y|$ 得 $|y|=\frac{8\sqrt3}{5}$。由方程求得
\[
x^2=17\left(1+\frac{y^2}{8}\right)=\frac{833}{25},
\quad P=\left(\frac{7\sqrt{17}}5,\frac{8\sqrt3}5\right)
\]
确在右支；其两焦半径为 $7+\sqrt{17}$、$7-\sqrt{17}$，满足差为 $2\sqrt{17}=2a$，且
\[
\cos\theta=\frac{(7+\sqrt{17})^2+(7-\sqrt{17})^2-10^2}
{2(7+\sqrt{17})(7-\sqrt{17})}=\frac12.
\]
故 $\theta=60^\circ$，数据自洽。

\medskip
\textbf{\textcolor{CExample}{例7（极限判断，不要误删顶点）}}
点 $P$ 在右支从顶点 $(a,0)$ 向远处移动时，面积 $S=c|y|$ 从 $0$ 逐渐增大；夹角 $\theta$ 则从 $\pi$ 逐渐趋向 $0$。
在顶点，两条焦半径指向相反方向，$\theta=\pi$、面积 $S=0$；虽然不是非退化三角形，但公式\emph{按连续延拓}给出 $b^2\cot(\pi/2)=0$。不能说顶点处夹角是 $0$ 或完全不存在。

\end{examplebox}

\clearpage
\begin{examplebox}
\textbf{\textcolor{CExample}{例8（旁切圆零余弦定理：三边 $6,4,4$）}}
取 $x^2-y^2/3=1$，$F_1(-2,0)$、$F_2(2,0)$，
$P=\bigl(\tfrac52,\tfrac{3\sqrt7}2\bigr)$。
代入方程：$25/4-(63/4)/3=1$，故 $P$ 确在双曲线上。
\[
 PF_1=6,\qquad PF_2=4,\qquad F_1F_2=4;
 \qquad a=1,\ c=2,\ b^2=3.
\]
\textbf{画出 $F_1$ 对应的旁切圆：} 令 $T$ 在线段 $PF_2$，$U$ 在 $F_1P$ 越过 $P$ 的延长线上，
$V$ 在 $F_1F_2$ 越过 $F_2$ 的延长线上。由等切线长，
\[
 PT=PU=\frac{4+4-6}{2}=\boxed{1},\qquad
 F_2T=F_2V=4-1=3.
\]
圆心到三个切点的距离都为半径 $r$，所以
\[
 S=\frac r2(6+4-4)=3r,
 \quad\text{且}\quad
 \frac r{PT}=\cot\frac\theta2.
\]
无需计算 $\cos\theta$，直接得到
\[
 S=3\cot\frac\theta2=b^2\cot\frac\theta2.
\]
\textbf{用底高校验：} 由给定坐标 $S=c|y|=2\cdot\tfrac{3\sqrt7}{2}=3\sqrt7$，
因 $S=3r$，可得 $r=\sqrt7$，相应 $\cot(\theta/2)=\sqrt7$，完全一致。
这一例的\textbf{切点与半径标注}见方法五的精确图。
\end{examplebox}
